\chapter{Ejercicios}

\section{Ejercicios del capítulo 1}

\begin{ejer}\label{ejer:identidad-auxiliar-disco-unidad}
    Sean $\alpha, \beta \in \mathbb{D}$. Demuestra la siguiente identidad:
    \[
    1 - \abs{\frac{\alpha - \beta}{1 - \overline{\beta} \alpha}}^2 = \frac{(1 - \abs{\alpha}^2)(1 - \abs{\beta}^2)}{\abs{1 - \overline{\beta} \alpha}^2}.
    \]
\end{ejer}
\begin{sol}
    Partiendo del lado izquierdo, tenemos que
    \[
    1 - \abs{\frac{\alpha - \beta}{1 - \bar{\beta} \alpha}}^2 = \frac{\abs{1 - \bar{\beta} \alpha}^2 - \abs{\alpha - \beta}^2}{\abs{1 - \bar{\beta} \alpha}^2}.
    \]
    En el numerador, desarrollando ambos términos, obtenemos
    \[
    \begin{aligned}
        \abs{1 - \bar{\beta} \alpha}^2 - \abs{\alpha - \beta}^2 &= \left(1 - \bar{\beta} \alpha\right)\left(1 - \beta \bar{\alpha}\right) - (\alpha - \beta)(\bar{\alpha} - \bar{\beta}) \\
        &= 1 - \cancel{\bar{\alpha} \beta} - \cancel{\alpha \bar{\beta}} + \abs{\alpha}^2 \abs{\beta}^2 - \left(\abs{\alpha}^2 - \cancel{\alpha \bar{\beta}} - \cancel{\bar{\alpha} \beta} + \abs{\beta}^2\right) \\
        &= 1 - \abs{\alpha}^2 - \abs{\beta}^2 + \abs{\alpha}^2 \abs{\beta}^2 = (1 - \abs{\alpha}^2)(1 - \abs{\beta}^2).
    \end{aligned}
    \]
    Por tanto, se sigue la identidad deseada.\hfill \qedsymbol
\end{sol}

\begin{ejer}\label{ejer:desigualdad-schwarz-pick}
    Sea $f \in \exhyperref[defn:clase-schur]{clase-schur}{\mathcal{S}}$. Demuestra que
    \[
    \forall z, w \in \mathbb{D} : \abs{\frac{f(z) - f(w)}{z-w}}^2 \leq \frac{1 - \abs{f(z)}^2}{1 - \abs{z}^2} \cdot \frac{1 - \abs{f(w)}^2}{1 - \abs{w}^2}.
    \]
\end{ejer}
\begin{sol}
    Como $f \in \mathcal{S}$, por el \exhyperref[teo:schwarz-pick]{teo-schwarz-pick}{teorema de Schwarz-Pick}, tenemos que
    \[
    \abs{\frac{f(z) - f(w)}{1 - \overline{f(w)} f(z)}} \leq \abs{\frac{z - w}{1 - \bar{w} z}} \implies 1 - \abs{\frac{z - w}{1 - \bar{w} z}}^2 \leq 1 - \abs{\frac{f(z) - f(w)}{1 - \overline{f(w)} f(z)}}^2. 
    \]

    Por el \excref[ejer:identidad-auxiliar-disco-unidad]{ejer-identidad-auxiliar-disco-unidad}, esta desigualdad se traduce a
    \[
    \begin{aligned}
        \frac{(1 - \abs{z}^2)(1 - \abs{w}^2)}{\abs{1 - \bar{w} z}^2} &\leq \frac{(1 - \abs{f(z)}^2)(1 - \abs{f(w)}^2)}{\abs{1 - \overline{f(w)} f(z)}^2} \\
        \implies \abs{\frac{1 - \bar{f(w)} f(z)}{1 - \bar{w} z}}^2 &\leq \frac{1 - \abs{f(z)}^2}{1 - \abs{z}^2} \cdot \frac{1 - \abs{f(w)}^2}{1 - \abs{w}^2}.
    \end{aligned}
    \]

    Finalmente, observamos que:
    \[
    \begin{aligned}
        \abs{\frac{f(z) - f(w)}{z - w}}^2 &= \abs{\frac{f(z) - f(w)}{1 - \overline{f(w)} f(z)}}^2 \cdot \abs{\frac{1 - \overline{f(w)} f(z)}{1 - \bar{w} z}}^2 \cdot \abs{\frac{1 - \bar{w} z}{z - w}}^2 \\
        &\leq 1 \cdot \frac{1 - \abs{f(z)}^2}{1 - \abs{z}^2} \cdot \frac{1 - \abs{f(w)}^2}{1 - \abs{w}^2} \cdot 1 = \frac{1 - \abs{f(z)}^2}{1 - \abs{z}^2} \cdot \frac{1 - \abs{f(w)}^2}{1 - \abs{w}^2}.
    \end{aligned}
    \]

    \vspace{-\baselineskip}
    \vspace{-\belowdisplayskip}
    \hfill \qedsymbol
\end{sol}

\section{Ejercicios del capítulo 2}

\begin{ejer}
    Demuestra que la métrica de Poincaré $\wp$ satisface que
    \[
    \forall z, w \in \mathbb{D} : \wp(z, w) = 2 \tanh^{-1} \abs{\frac{z - w}{1 - \overline{w}z}}.
    \]
\end{ejer}
\begin{sol}
    Recordemos la definición de la tangente hiperbólica:
    \[
    \tanh(z) = \frac{\sinh(x)}{\cosh(x)} = \frac{e^x - e^{-x}}{e^x + e^{-x}} = \frac{e^{2x} - 1}{e^{2x} + 1}.
    \]
    Entonces, despejando, tenemos que $e^{2x} \left(1 - \tanh(x)\right) = 1 + \tanh(x)$, luego
    \[
    e^{2x} = \frac{1 + \tanh(x)}{1 - \tanh(x)} \implies \tanh^{-1}(y) = \frac{1}{2} \, \log \frac{1 + y}{1 - y}.
    \]
    En particular, para $y = \varrho (z, w)$ con $z, w \in \mathbb{D}$, tenemos que
    \[
    \tanh^{-1} \left(\varrho(z, w)\right) = \frac{1}{2} \, \log \frac{1 + \varrho(z, w)}{1 - \varrho(z, w)} = \frac{1}{2} \, \wp(z, w)
    \]
    por definición de la métrica de Poincaré. Por tanto, como $\varrho(z, w) = \abs{\frac{z - w}{1 - \overline{w} z}}$, se tiene que
    \[
    \implies \wp(z, w) = 2 \, \tanh^{-1} \left(\varrho(z, w)\right) = 2 \, \tanh^{-1} \abs{\frac{z - w}{1 - \overline{w} z}}.
    \]

    \vspace{-\baselineskip}
    \vspace{-\belowdisplayskip}
    \hfill\qedsymbol
\end{sol}

\begin{ejer}
    Sea $f \colon \Omega_1 \to \Omega_2$ una aplicación holomorfa entre dos dominios $\Omega_1, \Omega_2 \subset \C$ dotados de densidades métricas $\mu_1, \mu_2$ respectivamente. Demuestra que
    \[
    f \tex{ es una isometría entre } (\Omega_1, d_{\mu_1}) \tex{ y } (\Omega_2, d_{\mu_2}) \iff \forall z \in \Omega_1 : (f^* \mu_2)(z) = \mu_1(z).
    \]
\end{ejer}
% \begin{sol}
%     Veamos ambas implicaciones por separado:
%     \begin{itemize}
%         \item[$\boxed{\Rightarrow}$]
%         \item[$\boxed{\Leftarrow}$] 
%     \end{itemize}
% \end{sol}
%DEMOSTRACIÓN

\begin{ejer}%EJERCICIO 
    Demuestra que la curvatura de la métrica esférica es constante e igual a $+1$.
\end{ejer}
% \begin{sol}

% \end{sol}

%COMPLETAR mirar Geometry: Euclid and Beyond.


